\SourcePage{318}
\SGAsection{5}{旗概形}
设 $S$ 是一个\Term{scheme}{概形}，$E$ 是\Term{rank}{秩}为 $r$ 的\Term{locally-free}{局部自由} $\mathscr O_S$-\Term{module}{模}，$(p_i)_{1\le i\le m}$ 是 $m$ 个正整数，满足
\[
 p_1+p_2+\cdots+p_m=r.
\]
回顾一下，$E$ 的一个 $(p_1,\ldots,p_m)$ 型\Term{flag}{旗}，是指 $E$ 的一个\Term{finite}{有限}\Term{filtration}{滤过}
\[
 0=E_m\subset E_{m-1}\subset\cdots\subset E_0=E,
\]
使得对每个整数 $i\in[1,m]$，\Term{quotient}{商} $E_{i-1}/E_i$ 都是\Term{rank}{秩}为 $p_i$ 的\Term{locally-free}{局部自由}\Term{module}{模}。

定义\Term{functor}{函子} $(\operatorname{Sch}/S)^\circ\to\operatorname{Ens}$：把每个 $S$-\Term{scheme}{概形} $u:T\to S$ 映为 $u^*(E)$ 的所有 $(p_1,\ldots,p_m)$ 型\Term{flag}{旗}组成的\Term{set}{集合}，把 $S$-\Term{scheme}{概形}之间的每个\Term{morphism}{态射}映为\Term{flag}{旗}上意义显然的\Term{inverse-image}{逆像}运算。可以证明（EGA I，第 2 版），这个\Term{functor}{函子}是\Term{representable}{可表的}；它的任一代表称为 $E$ 的 $(p_1,\ldots,p_m)$ 型\Term{flag-scheme}{旗 $S$-概形}。

以下用 $f:X\to S$ 表示 $E$ 的 $(p_1,\ldots,p_m)$ 型\Term{flag-scheme}{旗 $S$-概形}。此外，把 $E$ 的 $(1,1,\ldots,1)$ 型\Term{flag-scheme}{旗概形}记为 $D$；在不会混淆时，径称其为 $E$ 的\Term{flag-scheme}{旗概形}。

由定义显然可见，$\mathscr O_X$-\Term{module}{模} $f^*(E)$ 有一个\Term{finite}{有限}的\Term{canonical}{典范}\Term{filtration}{滤过}
\[
 0=F_m\subset F_{m-1}\subset\cdots\subset F_0=f^*(E),
\]
其中\Term{quotient}{商} $E_i=F_{i-1}/F_i$（$1\le i\le m$）是\Term{rank}{秩}为 $p_i$ 的\Term{locally-free}{局部自由}\Term{module}{模}。对每个整数 $i\in[1,m]$，以
\[
 g_i:D_i\longrightarrow X
\]
表示 $E_i$ 的\Term{flag-scheme}{旗 $X$-概形}，并以 $g:D_1\times_X\cdots\times_XD_m\to X$ 表示这些 $X$-\Term{scheme}{概形} $D_i$ 的\Term{fiber-product}{纤维积}。显然，\Term{composite-morphism}{复合态射}
\SourcePage{319}
\[
 h=f\circ g:\prod_i{}_{\!X}D_i\longrightarrow S\tag{5.1}
\]
与 $E$ 的\Term{flag-scheme}{旗 $S$-概形} $S$-\Term{isomorphism}{同构}。因此，在相差一个 $S$-\Term{isomorphism}{同构}的意义下，可写成 $D=\prod_i{}_{\!X}D_i$。

\begin{SGAresult}{命题}{5.2}
在 $D^b(S,A)$ 中有\Term{isomorphism}{同构}
\[
 \mathbf Rf_*(A)\simeq\bigoplus_iR^{2i}f_*(A)[-2i].
\]
\end{SGAresult}
\begin{proof}
由于\Term{morphism}{态射} $h$ 是\Term{finite}{有限}个如下类型的\Term{morphism}{态射}的复合：
\[
 (Q)\qquad \mathbf P_T(F)\longrightarrow F
 \qquad(F\text{ 为局部自由的 }\mathscr O_T\text{-模}),
\]
由 (2.2.1) 得
\[
 \mathbf Rh_*(A)\simeq\bigoplus_jR^jh_*(A)[-j].
\]
同理也有
\[
 \mathbf Rg_*(A)\simeq\bigoplus_jR^jg_*(A)[-j].
\]
利用 [6] 2.16，得到\Term{isomorphism}{同构}
\[
 \mathbf Rf_*(g_*A)\simeq\bigoplus_jR^jf_*(g_*A)[-j].
\]
由于 $g$ 是\Term{smooth}{光滑}的，且其\Term{fiber}{纤维}\Term{geometrically-connected}{几何连通}，有 $A_X\simeq g_*(A_D)$，所以只需证明
\[
 R^{2i+1}f_*(A)=0\qquad(i\in\mathbf N).\tag{5.2.1}
\]
由\Term{fiber-functor}{纤维函子}系的\Term{conservative}{保守性}（SGA 4 第 VIII 讲 3.5），立即归结为 $S$ 是\Term{separably-closed-field}{可分闭域} $k$ 的谱的情形。这时，\Term{morphism}{态射} $g$ 是若干\Term{projective-bundle}{射影丛}的复合，由 (2.2.4)，\Term{inverse-image}{逆像}\Term{morphism}{态射}
\[
 H^{2i+1}(X,A_X)\longrightarrow H^{2i+1}(D,A_D)
\]
是\Term{monomorphism}{单态射}，故可假设 $X=D$。在这种情形下，
\SourcePage{320}
\Term{projection}{投影}\Term{morphism}{态射} $D\to k$ 本身也是若干\Term{projective-bundle}{射影丛}的复合，因而对 (2.2.4) 作归纳便得断言。
\end{proof}

\begin{SGAresult}{推论}{5.3}
对任意\Term{scheme}{概形}\Term{morphism}{态射} $u:S\to T$，Leray \Term{spectral-sequence}{谱序列}
\[
 E_2^{pq}=R^pu_*R^qf_*(A)\Longrightarrow R^{p+q}(u\circ f)_*(A)
\]
是\Term{degenerate-spectral-sequence}{退化的}。
\end{SGAresult}
\begin{proof}
见 [6] 1.2。
\end{proof}

现在进一步确定 $X$ 的\Term{cohomology}{上同调}的\Term{multiplication}{乘法}结构，先考察一个特例。沿用上述记号，借助\Term{inverse-image}{逆像} $f^*$ 把 $A^\bullet(S)$ 映入 $A^\bullet(X)$，并把每个如下指标对 $(i,j_i)$ 对应的元素 $T_{ij_i}$ 映为第 $j_i$ 个\Term{chern-class}{陈类} $c_{j_i}(E_i)$，便得到\Term{ring-homomorphism}{环同态}
\[
 \alpha:A^\bullet(S)[T_{ij_i}]_{1\le i\le m,\,1\le j_i\le p_m}
 \longrightarrow A^\bullet(X).
\]
对每个整数 $p$，由 (3.4 (iii)) 有
\[
 c_p(f^*E)=\sum_{i_1+i_2+\cdots+i_m=p}
 c_{i_1}(E_1)\cdots c_{i_m}(E_m).
\]
因此，由\Term{chern-class}{陈类}的\Term{functoriality}{函子性}，这个\Term{morphism}{态射}在由所有元素
\[
 c_p(E)-\sum_{i_1+\cdots+i_m=p}T_{1,i_1}\cdots T_{m,i_m}
\]
（$p$ 遍历各整数）生成的\Term{ideal}{理想} $J_X$ 上为零。于是得到\Term{ring-homomorphism}{环同态}
\[
 \varphi_X:\bigl(A^\bullet(S)[T_{ij_i}]_{(i,j_i)}\bigr)/J_X
 \longrightarrow A^\bullet(X).
\]
\begin{SGAresult}{命题}{5.4}
上述\Term{ring-homomorphism}{环同态} $\varphi_X$ 是\Term{isomorphism}{同构}。$A^\bullet(S)$-\Term{module}{模} $A^\bullet(X)$ 是\Term{rank}{秩}为
\[
 \frac{r!}{p_1!p_2!\cdots p_m!}
\]
的\Term{free-module}{自由模}。
\end{SGAresult}
\SourcePage{321}
\begin{proof}
遵循 Grothendieck 在 Chevalley 讨论班 1958 年《Chow 环及其应用》第 IV 讲中的论述。需要如下\Term{lemma}{引理}（参见该文第 19 页\Term{lemma}{引理} 1）。
\end{proof}
\begin{SGAresult}{\Term{lemma}{引理}}{5.4.1}
a) 设 $R$ 是\Term{commutative-ring}{交换环}，$(a_i)_{1\le i\le r}$ 是 $R$ 的 $r$ 个元素。记 $\sigma_i$ 为元素 $(U_i)_{1\le i\le r}$ 的第 $i$ 个\Term{elementary-symmetric-function}{初等对称函数}，$I$ 为\Term{polynomial-ring}{多项式环} $R[U_1,U_2,\ldots,U_r]$ 中由元素 $\sigma_i-a_i$ 生成的\Term{ideal}{理想}，则 $R$-\Term{algebra}{代数}
\[
 R[U_1,U_2,\ldots,U_r]/I
\]
是\Term{rank}{秩}为 $r!$ 的\Term{free-module}{自由 $R$-模}。

b) 设 $r,p_1,p_2,\ldots,p_m$ 是正整数，满足 $p_1+\cdots+p_m=r$。将 $T_{i,j_i}$ 映为元素
\[
 (U_{p_1+\cdots+p_{i-1}+s})_{1\le s\le p_i}
\]
的第 $j_i$ 个\Term{elementary-symmetric-function}{初等对称函数}所得到的\Term{polynomial-ring}{多项式环}\Term{morphism}{态射}
\[
 V=R[T_{ij_i}]_{1\le i\le m,\,1\le j_i\le p_m}
 \longrightarrow R[U_i]_{1\le i\le r}=W
\]
是\Term{injective}{单射}，且使 $W$ 成为\Term{rank}{秩}为 $p_1!\cdots p_m!$ 的\Term{free-module}{自由 $V$-模}。

特别地，把 $T_i$ 映为各 $U_i$ 的第 $i$ 个\Term{elementary-symmetric-function}{初等对称函数}所得到的 $R$-\Term{algebra}{代数}\Term{morphism}{态射}
\[
 R[T_1,\ldots,T_r]\longrightarrow R[U_1,\ldots,U_r]
\]
是\Term{injective}{单射}，且使右端成为左端上的\Term{rank}{秩}为 $r!$ 的\Term{free-module}{自由模}。
\end{SGAresult}
\begin{proof}[命题 5.4 的证明（续）]
说明 (5.4.1) 如何\Term{pushout}{推出} (5.4)。先设 $X$ 是 $E$ 的\Term{flag-scheme}{旗概形} $D$，并为简便置 $T_{i1}=U_i$。
\SourcePage{322}
由于\Term{canonical-projection}{典范投影} $D\to S$ 是 $r-1$ 个形如
\[
 \mathbf P_{Y_i}(M_i)\longrightarrow Y_i\qquad(1\le i\le r-1)
\]
的\Term{morphism}{态射}的复合，其中 $M_i$ 是\Term{rank}{秩}为 $r-i+1$ 的\Term{locally-free}{局部自由} $\mathscr O_{Y_i}$-\Term{module}{模}，由 (2.2.6) 得 $A^\bullet(D)$ 是\Term{rank}{秩}为 $r!$ 的\Term{free-module}{自由 $A^\bullet(S)$-模}，且\Term{morphism}{态射} $\varphi_D$ 是\Term{surjective}{满射}。但\Term{lemma}{引理} 5.4.1 a) 表明，$\varphi_D$ 的源也是\Term{rank}{秩}为 $r!$ 的\Term{free-module}{自由 $A^\bullet(S)$-模}，这就证明了 $X=D$ 的情形。在一般情形下，显然下图\Term{commutative-diagram}{交换}：
\[
\begin{tikzcd}[column sep=large,row sep=large]
 A^\bullet(X)\arrow[r,"g^*"]&A^\bullet(D)\\
 C=A^\bullet(S)[T_{i,j_i}]/J_X\arrow[u,"\varphi_X"]\arrow[r,"\gamma"]
 &A^\bullet(S)[U_k]/J_D\arrow[u,"\varphi_D"']
\end{tikzcd}
\]
其中映射 $\gamma$ 是由 $A^\bullet(S)$-\Term{algebra}{代数}\Term{morphism}{态射}
\[
 A^\bullet(S)[T_{i,j_i}]\longrightarrow A^\bullet[U_k]
\]
取\Term{quotient}{商}得到的，而此\Term{morphism}{态射}将 $T_{i,j_i}$ 映为元素
\[
 (U_{p_1+\cdots+p_{i-1}+s})_{1\le s\le p_i}
\]
的第 $j_i$ 个\Term{elementary-symmetric-function}{初等对称函数}。由于 $\varphi_D$ 是\Term{isomorphism}{同构}，由 5.4.1 b)，$A^\bullet(D)$ 是\Term{rank}{秩}为 $p_1!\cdots p_m!$ 的\Term{free-module}{自由 $C$-模}。另一方面，将 $g$ 表为若干\Term{flag-scheme}{旗 $X$-概形}的\Term{fiber-product}{纤维积}，可见 $A^\bullet(D)$ 也是\Term{rank}{秩}为 $p_1!\cdots p_m!$ 的\Term{free-module}{自由 $A^\bullet(X)$-模}。由此立即得到 $\varphi_X$ 是\Term{isomorphism}{同构}。
\end{proof}

在下一命题中，约定对每个\Term{scheme}{概形} $Y$ 和每个 $A_Y$-\Term{module}{模} $M$，置
\[
 Q(Y,M)=\bigoplus_{n,i}H^n(Y,M\otimes\mu_Y^{\otimes i}).
\]
\SourcePage{323}
\Term{cup-product}{杯积}使 $A$-\Term{module}{模} $Q(Y,M)$ \Term{canonical}{典范}地具有 $A^\bullet(Y)$-\Term{module}{模}结构。若 $u:Z\to Y$ 是\Term{scheme}{概形}\Term{morphism}{态射}，则\Term{inverse-image}{逆像}\Term{morphism}{态射}
\[
 u^*:Q(Y,M)\longrightarrow Q(Z,u^*M)
\]
与 $Q(Y,M)$ 的 $A^\bullet(Y)$-\Term{module}{模}结构及 $Q(Z,u^*M)$ 的 $A^\bullet(Z)$-\Term{module}{模}结构相容。由此得到一个称为\Term{canonical}{典范}\Term{morphism}{态射}的 $A^\bullet(Z)$-\Term{module}{模}\Term{morphism}{态射}
\[
 A^\bullet(Z)\otimes_{A^\bullet(Y)}Q(Y,M)\longrightarrow Q(Z,u^*M).
\]
\begin{SGAresult}{命题}{5.5}
设 $S$ 是\Term{scheme}{概形}，$E$ 是\Term{rank}{秩}为 $r$ 的\Term{locally-free}{局部自由} $\mathscr O_S$-\Term{module}{模}。以 $f:X\to S$ 表示 $E$ 的 $(p_1,\ldots,p_m)$ 型\Term{flag-scheme}{旗 $S$-概形}。则对每个 $A_S$-\Term{module}{模} $L$，\Term{canonical}{典范}\Term{morphism}{态射}
\[
 \alpha:A^\bullet(X)\otimes_{A^\bullet(S)}Q(S,L)\longrightarrow Q(X,f^*L)
\]
是\Term{isomorphism}{同构}。
\end{SGAresult}
\begin{proof}
当 $X=\mathbf P_S(E)$ 时，这立即由 (2.2.4) 和 (2.2.6) 得到。由此可得 $X=D$ 时也成立。在一般情形下，由于\Term{canonical}{典范}\Term{morphism}{态射} $g:D\to X$ 和 $h:D\to S$ 是若干\Term{projective-bundle}{射影丛}的复合，\Term{canonical}{典范}\Term{morphism}{态射}
\[
 A^\bullet(D)\otimes_{A^\bullet(S)}Q(S,L)\longrightarrow Q(D,h^*L)
\]
及
\[
 A^\bullet(D)\otimes_{A^\bullet(X)}Q(X,g^*L)\longrightarrow Q(D,g^*L)
\]
都是\Term{isomorphism}{同构}。所以 $\operatorname{id}_{A^\bullet(D)}\otimes\alpha$ 是\Term{isomorphism}{同构}；由于 $A^\bullet(D)$ 是 $A^\bullet(X)$ 上的\Term{free-module}{自由模}，结论成立。
\end{proof}

现在给出 (5.4) 和 (5.5) 的相对版本。沿用上述记号，设 $u:S\to T$ 是\Term{scheme}{概形}\Term{morphism}{态射}。置
\[
 A^\bullet(X/T)=\bigoplus_iR^{2i}(u\circ f)_*(A_X),
\]
同样置 $A^\bullet(S/T)=\bigoplus_iR^{2i}u_*(A_S)$。最后，对每个 $A_S$-\Term{module}{模} $L$，置
\SourcePage{324}
\[
 Q(S/T,L)=\bigoplus_{n,i}R^{2i}u_*(L\otimes\mu_S^{\otimes i})
\]
及
\[
 Q(X/T,f^*L)=\bigoplus_{n,i}R^n(u\circ f)_*(f^*L\otimes\mu_X^{\otimes i}).
\]
对每对满足 $1\le i\le m$、$1\le j\le p_i$ 的整数 $(i,j)$，\Term{chern-class}{陈类} $c_j(E_i)$ 在\Term{canonical}{典范}\Term{morphism}{态射}
\[
 H^{2j}(X,\mu_X^{\otimes j})\longrightarrow
 H^0(T,R^{2j}(u\circ f)_*(\mu_X^{\otimes j}))
\]
下的像，定义了一个 $A_T$-\Term{morphism}{态射}
\[
 c_{ij}:A_T\longrightarrow R^{2j}(u\circ f)_*(\mu_X^{\otimes j}).
\]
这些 $c_{ij}$ 以及\Term{inverse-image}{逆像} $f^*$ 定义了 $A^\bullet(S/T)$-\Term{algebra}{代数}\Term{morphism}{态射}
\[
 \theta_1:A^\bullet(S/T)[T_{i,j_i}]_{(i,j_i)}\longrightarrow A^\bullet(X/T).
\]
以 $c_i$ 表示由\Term{chern-class}{陈类} $c_i(E)$ 定义的 $A^\bullet(S/T)$ 的\Term{section}{截面}。与 (5.4) 一样可见，$\theta_1$ 在\Term{local-scheme}{局部}由关系
\[
 \sum_{i_1+\cdots+i_m=p}T_{1,i_1}\cdots T_{m,i_m}=c_p\qquad(p\in\mathbf N)
\]
生成的\Term{ideal}{理想} $J$ 上为零。于是得到 $A^\bullet(S/T)$-\Term{algebra}{代数}\Term{morphism}{态射}
\[
 \theta:A^\bullet(S/T)[T_{i,j_i}]_{(i,j_i)}/J\longrightarrow A^\bullet(X/T).
\]
\begin{SGAresult}{命题}{5.6}
沿用上述记号。

a) $A^\bullet(S/T)$-\Term{algebra}{代数}\Term{morphism}{态射}
\[
 \theta:A^\bullet(S/T)[T_{i,j_i}]_{(i,j_i)}/J\longrightarrow A^\bullet(X/T)
\]
是\Term{isomorphism}{同构}。

b) 对每个 $A_S$-\Term{module}{模} $L$，\Term{canonical}{典范}的 $A^\bullet(X/T)$-\Term{module}{模}\Term{morphism}{态射}
\SourcePage{325}
\[
 A^\bullet(X/T)\otimes_{A^\bullet(S/T)}Q(S/T,L)\longrightarrow Q(X/T,f^*L)
\]
是\Term{isomorphism}{同构}。
\end{SGAresult}
\begin{proof}
分别将 (5.4) 和 (5.5) 作\Term{sheafification}{层化}即可。
\end{proof}

\SourcePage{326}
\SGAsection{6}{群概形}
先讨论由\Term{torsor}{挠子}\Term{twisting}{扭曲}的\Term{object}{对象}的一条一般性质。

\begin{SGAresult}{命题}{6.1}
设 $S$ 是\Term{scheme}{概形}，$G$ 是\Term{quasi-compact}{拟紧}、\Term{smooth}{光滑}且\Term{fiber}{纤维}\Term{connected}{连通}的 $S$-\Term{group-scheme}{群概形}，$X$ 是带有 $G$ 的\Term{left-action}{左作用}的 $S$-\Term{scheme}{概形}。若 $f:E\to S$ 是 $G$ 下的一个右 $S$-\Term{torsor}{挠子}，以 $X_E$ 表示借助 $E$ 及 $G$ 在 $X$ 上的\Term{action}{作用}，对 $X$ 作\Term{twisting}{扭曲}得到的 $S$-\Term{scheme}{概形}。以 $p:X\to S$ 和 $q:X_E\to S$ 表示\Term{canonical-projection}{典范投影}。则对每个 $A_S$-\Term{module}{模} $M$ 和每个整数 $i$，$A_S$-\Term{module}{模} $R^iq_*(q^*M)$ 与 $R^ip_*(p^*M)$ \Term{canonical}{典范}\Term{isomorphism}{同构}。
\end{SGAresult}
\begin{proof}
考虑下列两图：
\[
\begin{tikzcd}[column sep=large,row sep=large]
 E\times_SX\arrow[r,"\varphi"]\arrow[d,"\operatorname{pr}_1"']&X_E\arrow[d,"q"]\\
 E\arrow[r,"f"']&S
\end{tikzcd}
\qquad(D_1)
\qquad
\begin{tikzcd}[column sep=large,row sep=large]
 E\times_SX\arrow[r,"\operatorname{pr}_2"]\arrow[d,"\operatorname{pr}_1"']&X\arrow[d,"p"]\\
 E\arrow[r,"f"']&S
\end{tikzcd}
\qquad(D_2).
\]
已知图 $(D_1)$ 是\Term{cartesian-diagram}{笛卡尔图}，其中 $\varphi$ 是取\Term{quotient}{商}的\Term{canonical-map}{典范映射}。由于\Term{morphism}{态射} $f$ 是\Term{smooth}{光滑}的，\Term{smooth-base-change}{光滑基变换定理}给出\Term{isomorphism}{同构}
\[
 u:f^*R^iq_*(q^*M)\xrightarrow{\sim}
 R^i(\operatorname{pr}_1)_*(\varphi^*q^*M).
\]
同样，由\Term{cartesian-diagram}{笛卡尔图} $(D_2)$ 得到\Term{isomorphism}{同构}
\[
 v:f^*R^ip_*(p^*M)\xrightarrow{\sim}
 R^i(\operatorname{pr}_1)_*((\operatorname{pr}_2)^*p^*M).
\]
但 $p\circ\operatorname{pr}_2=q\circ\varphi$，所以将 $\varphi^*\circ q^*$ 与 $(\operatorname{pr}_2)^*\circ p^*$ 按显然的方式等同，便定义了\Term{isomorphism}{同构}
\[
 v^{-1}\circ u:f^*R^iq_*(q^*M)\xrightarrow{\sim}f^*R^ip_*(p^*M).
\]
\Term{morphism}{态射} $f$ 的\Term{geometric-fiber}{几何纤维}\Term{nonempty}{非空}且\Term{connected}{连通}，
\SourcePage{327}
由 SGA 4 第 XII 讲 6.5 (i)，\Term{adjunction-morphism}{伴随态射} $\operatorname{adj}:\operatorname{id}\to f_*f^*$ 是\Term{isomorphism}{同构}，因此
\[
 (\operatorname{adj})^{-1}\circ f_*(v^{-1}\circ u)\circ(\operatorname{adj})
\]
就是所求的\Term{isomorphism}{同构}。
\end{proof}

\begin{SGAresult}{命题}{6.2}
（M. Raynaud 女士 [5]。）设 $S$ 是\Term{scheme}{概形}，$p:G\to S$ 是 $S$ 上的\Term{reductive-group-scheme}{约化群概形}（SGA 3 第 XIX 讲）。对每个\Term{locally-constant}{局部常值}\Term{constructible}{可构造} $A_S$-\Term{module}{模} $M$ 和每个整数 $i$，$A_S$-\Term{module}{模} $R^ip_*(p^*M)$ 都是\Term{locally-constant}{局部常值}\Term{constructible}{可构造}的。此外，对每个\Term{morphism}{态射} $u:S'\to S$ 及其定义的\Term{cartesian-diagram}{笛卡尔图}
\[
\begin{tikzcd}[column sep=large,row sep=large]
 G'\arrow[r,"u'"]\arrow[d,"p'"']&G\arrow[d,"p"]\\
 S'\arrow[r,"u"']&S,
\end{tikzcd}
\]
\Term{base-change-morphism}{基变换态射}
\[
 u^*R^ip_*(p^*M)\longrightarrow R^ip'_*(p\circ u')^*(M)
\]
是\Term{isomorphism}{同构}。
\end{SGAresult}
\begin{proof}
命题中的两个断言对 $S$ 的\Term{etale-topology}{平展拓扑}都是\Term{local-scheme}{局部}的。由 SGA 3 第 XII 讲 2.3，可假设 $G$ 是\Term{split}{分裂}的。取 $G$ 的一个如 SGA 3 第 XXII 讲 5.1.1 所示类型的\Term{borel-subgroup}{Borel 子群} $B$。由 SGA 3 第 XXII 讲 5.8.5，\Term{quotient}{商} $G/B$ 存在。设 $q:G\to G/B$ 和 $r:G/B\to S$ 是\Term{canonical}{典范}\Term{morphism}{态射}。由于\Term{morphism}{态射} $r$ 是\Term{proper}{固有}且\Term{smooth}{光滑}的，\Term{proper-base-change}{固有基变换定理}对它成立，且\Term{functor}{函子} $R^ir_*$ 将\Term{locally-constant}{局部常值}\Term{constructible}{可构造}的阿贝尔\Term{sheaf}{层}变为同类型的\Term{sheaf}{层}（SGA 4 第 XVI 讲 2.2）。另一方面，由 SGA 4 第 XXII 讲 5.5.1，$S$-\Term{scheme}{概形} $B$ 与\Term{finite}{有限}个\Term{vector-bundle}{向量丛}及\Term{punctured-vector-bundle}{去零向量丛}的积\Term{isomorphism}{同构}。因此由 (1.1) 和 (1.2) 立即可知，若以 $t:B\to S$ 表示\Term{canonical}{典范}\Term{morphism}{态射}，
\SourcePage{328}
则对\Term{locally-constant}{局部常值}\Term{constructible}{可构造}的 $A_S$-\Term{module}{模} $M$，\Term{sheaf}{层} $R^it_*(t^*M)$ 都是\Term{locally-constant}{局部常值}\Term{constructible}{可构造}的，且在\Term{base-change}{基变换}下\Term{stable}{稳定}。现在取一个\Term{locally-constant}{局部常值}\Term{constructible}{可构造}的 $A_S$-\Term{module}{模} $M$。由上述讨论以及 (6.1)，\Term{sheaf}{层} $R^iq_*(p^*M)$ 都是\Term{locally-constant}{局部常值}\Term{constructible}{可构造}的。利用 Leray \Term{spectral-sequence}{谱序列}
\[
 R^ir_*(R^jq_*)(p^*M)\Longrightarrow R^{i+j}p_*(p^*M),
\]
可知\Term{sheaf}{层} $R^ip_*(p^*M)$ 也有同样的性质。最后，由 (6.1)，\Term{sheaf}{层} $R^iq_*(p^*M)$ 在\Term{base-change}{基变换}下\Term{stable}{稳定}；再对 $r$ 应用\Term{proper-base-change}{固有基变换定理}，便得\Term{sheaf}{层} $R^ip_*(p^*M)$ 也在\Term{base-change}{基变换}下\Term{stable}{稳定}。
\end{proof}

\textbf{注记 6.3。} 每个 $S$ 上的\Term{reductive-group}{约化群}，对\Term{etale-topology}{平展拓扑}\Term{local-scheme}{局部}而言，都是某个 $\operatorname{Spec}(\mathbf Z)$ 上的\Term{reductive-group-scheme}{约化群概形}沿\Term{canonical}{典范}\Term{morphism}{态射} $S\to\operatorname{Spec}(\mathbf Z)$ 的\Term{inverse-image}{逆像}（SGA 3 第 XXIII 讲 5.7 及第 XXV 讲 1.3）。因而 (6.2) 表明，\Term{algebraically-closed-field}{代数闭域}上的\Term{reductive-group-scheme}{约化群概形}的\Term{cohomology}{上同调}计算（系数为与\Term{characteristic}{特征}互素的\Term{torsion-coefficient}{挠系数}），可归结为复数域上同类型的\Term{reductive-group-scheme}{约化群概形}的\Term{cohomology}{上同调}计算这一超越问题。借助\Term{comparison-theorem}{比较定理}（SGA 4 第 XI 讲 4.4），这一问题的解可由\Term{algebra}{代数}\Term{topology}{拓扑}的方法得到。

\begin{SGAresult}{推论}{6.4}
设 $S$ 是\Term{scheme}{概形}，$p:G\to S$ 是 $S$-\Term{group-scheme}{群概形}，它是一个\Term{finite}{有限}\Term{etale}{平展} $S$-\Term{group-scheme}{群概形} $F$ 被一个 $S$ 上的\Term{reductive-group-scheme}{约化群概形} $H$ 的\Term{extension}{扩张}。则对每个\Term{locally-constant}{局部常值}\Term{constructible}{可构造} $A_S$-\Term{module}{模} $M$，(6.2) 的结论都成立。
\end{SGAresult}
\begin{proof}
作\Term{base-change}{基变换} $F\to S$ 后，可假设 $F$ 形如 $\underline F_S$，其中 $F$ 是一个通常的\Term{finite-group}{有限群}。这时有 $S$-\Term{scheme}{概形}的\Term{isomorphism}{同构}
\[
 G\simeq\coprod_{\gamma\in F}H,
\]
再应用 (6.2) 即得结论。
\end{proof}

\SourcePage{329}
\begin{SGAresult}{推论}{6.5}
设 $S$ 是\Term{scheme}{概形}，$G$ 是 $S$ 上的\Term{reductive-group-scheme}{约化群概形}，$H$ 是 $G$ 的一个\Term{reductive-group-scheme}{约化子 $S$-群概形}，并假设在\Term{faithfully-flat}{忠实平坦}\Term{quasi-compact}{拟紧}\Term{topology}{拓扑}下，$G$ 对 $H$ 的\Term{quotient-sheaf}{商层}是\Term{representable}{可表的}。以 $p:G/H\to S$ 表示\Term{canonical}{典范}\Term{morphism}{态射}，则对每个\Term{locally-constant}{局部常值}\Term{constructible}{可构造} $A_S$-\Term{module}{模} $M$，(6.2) 的结论都成立。
\end{SGAresult}
\begin{proof}
设 $f:G\to G/H$ 是\Term{canonical}{典范}\Term{morphism}{态射}。它是\Term{smooth}{光滑}的，因此对\Term{etale-sheaf}{平展层}是\Term{descent}{下降}\Term{morphism}{态射}（SGA 4 第 VIII 讲 9.1）。对每个整数 $r\ge0$，置
\[
 G_r=\underbrace{G\times_{G/H}G\times_{G/H}\cdots\times_{G/H}G}_{r+1\text{ 个因子}},
\]
由 SGA 3 第 V 讲 1.5，它\Term{canonical}{典范}地等同于
\[
 G\times_S\underbrace{H\times_S\cdots\times_SH}_{r\text{ 个因子}}.
\]
以 $f_r:G_r\to G/H$ 和 $g_r:G_r\to S$ 表示\Term{canonical}{典范}\Term{morphism}{态射}，特别地有 $f_0=f$。最后，对每个整数 $q\ge0$，以 $\mathscr H^q(M)$ 表示\Term{complex}{复形}
\[
 R^q(g_0)_*(f_0^*M)\rightrightarrows
 R^q(g_1)_*(f_1^*M)\mathrel{\substack{\longrightarrow\\[-.6ex]\longrightarrow\\[-.6ex]\longrightarrow}}
 \cdots\mathrel{\substack{\longrightarrow\\[-.6ex]\longrightarrow\\[-.6ex]\longrightarrow}}
 R^q(g_r)_*(f_r^*M)\cdots.
\]
根据（SGA \SGAsourcegap{原书在 SGA 后留空，\Term{descent-spectral-sequence}{下降谱序列}的讲、节引注缺失}），有\Term{descent}{下降}\Term{spectral-sequence}{谱序列}（或 Čech \Term{spectral-sequence}{谱序列}）
\[
 E(M):E_2^{pq}=H^p(\mathscr H^q(M))\Longrightarrow E^n=R^np_*(M).
\]
由于各\Term{scheme}{概形} $G_r$ 都是\Term{reductive-group-scheme}{约化群概形}，由 6.2，各项 $E_2^{pq}$ 都是\Term{locally-constant}{局部常值}\Term{constructible}{可构造}的，且与任意\Term{base-change}{基变换}相容。由此立即可知，各 $E^n$ 也有同样的性质。
\end{proof}

\SourcePage{330}
\SGAsection{7}{完全交}
本节固定一个\Term{separably-closed-field}{可分闭域} $k$。若 $\ell$ 是与 $\operatorname{car}(k)$ 不同的素数，对每个\Term{proper}{固有} $k$-\Term{scheme}{概形} $X$ 和每个整数 $i\ge0$，以 $b_\ell^i(X)$（在不会混淆时也写成 $b^i(X)$）表示 $X$ 对系数 $\mathbf Z_\ell$ 的第 $i$ 个\Term{betti-number}{Betti 数}。若 $X$ 在 $k$ 上\Term{projective}{射影}，则称\Term{vector-bundle}{向量丛} $\mathbf V(\mathscr L^\vee)$ 的一个与\Term{zero-section}{零截面}\Term{transverse}{横截}的\Term{section}{截面}的零点集为 $X$ 的\Term{hyperplane-section}{超平面截面}，其中 $\mathscr L$ 是一个\Term{ample}{丰沛}的\Term{invertible-module}{可逆 $\mathscr O_X$-模}。

\begin{SGAresult}{\Term{theorem}{定理}}{7.1}
（Lefschetz \Term{theorem}{定理}。）设 $X$ 是\Term{dimension}{维数}为 $n$ 的\Term{projective}{射影}、\Term{smooth}{光滑}、\Term{connected}{连通} $k$-\Term{scheme}{概形}，$Y$ 是 $X$ 的一个\Term{subscheme}{子概形}，为 $X$ 的 $d$ 个\Term{hyperplane-section}{超平面截面}之交，$L$ 是\Term{locally-constant}{局部常值}\Term{constructible}{可构造} $A_X$-\Term{module}{模}。则：

(i) \Term{inverse-image}{逆像}\Term{homomorphism}{同态}
\[
 H^i(X,L)\longrightarrow H^i(Y,L|Y)
\]
当 $i\le n-d-1$ 时是\Term{isomorphism}{同构}，当 $i=n-d$ 时是\Term{monomorphism}{单态射}。

(ii) 假设 $Y$ 在 $k$ 上\Term{smooth}{光滑}，且\Term{dimension}{维数}为 $n-d$。则 \Term{gysin-homomorphism}{Gysin 同态}
\[
 H^i(Y,(L|Y)\otimes\mu_Y^{\otimes-d})\longrightarrow H^{i+2d}(X,L)
\]
当 $i\ge n-d+1$ 时是\Term{isomorphism}{同构}，当 $i=n-d$ 时是\Term{epimorphism}{满态射}。
\end{SGAresult}
\begin{proof}
众所周知，\Term{hyperplane-section}{超平面截面}的补集是\Term{affine-open}{仿射开集}，因而 $U=X\setminus Y$ 是 $d$ 个\Term{affine-open}{仿射开集}的并。断言 (i) 立即由\Term{cohomology}{上同调}\Term{exact-sequence}{正合列}
\[
 \cdots\longrightarrow H_c^i(U,L|U)\longrightarrow H^i(X,L)
 \longrightarrow H^i(Y,L|Y)\xrightarrow{\partial}H_c^{i+1}(U,L|U)\longrightarrow\cdots
\]
以及下列\Term{lemma}{引理}得到。
\end{proof}
\begin{SGAresult}{\Term{lemma}{引理}}{7.1.1}
设 $U$ 是\Term{dimension}{维数}为 $n$ 的\Term{finite-type}{有限型} $k$-\Term{scheme}{概形}，它是 $d$ 个\Term{affine-open}{仿射开集}的并，$F$ 是一个 $A_U$-\Term{module}{模}。
\SourcePage{331}
a) 若 $F$ 是\Term{constructible}{可构造}的，则
\[
 H^i(U,F)=0\qquad\text{当 }i\ge n+d.
\]
b) 若 $F$ 是\Term{locally-constant}{局部常值}的，且 $U$ 在 $k$ 上\Term{smooth}{光滑}，则
\[
 H_c^i(U,F)=0\qquad\text{当 }i\le n-d.
\]
\end{SGAresult}
\begin{proof}
当 $d=1$ 时，断言 a) 就是 SGA 4 第 XIV 讲 3.2。在一般情形下，以 $\mathscr U$ 表示由 $d$ 个\Term{affine-open}{仿射开集}组成的 $U$ 的一个\Term{covering}{覆盖}，再用 Čech \Term{spectral-sequence}{谱序列}
\[
 H^p(\mathscr U,\mathscr H^q(F))\Longrightarrow H^{p+q}(U,F)
\]
即可得出 a)。

证明 b)。由于\Term{functor}{函子} $H_c^i(U,\cdot)$ 与\Term{filtered-inductive-limit}{滤过归纳极限}交换（SGA 4 第 VII 讲 3.3），可假设 $F$ 是\Term{constructible}{可构造}的；这时利用\Term{poincare-duality}{Poincaré 对偶}（SGA 4 第 XVIII 讲），断言由 a) 得到。
\end{proof}
\begin{proof}[\Term{theorem}{定理} 7.1 的证明（续）]
最后，考虑到 \Term{gysin-homomorphism}{Gysin 态射}是通过\Term{canonical}{典范}\Term{isomorphism}{同构}
\[
 H^{i-2d}(Y,(L|Y)\otimes\mu_Y^{\otimes-d})\xrightarrow{\sim}H_Y^i(X,L)\tag{7.1.2}
\]
定义的，断言 (ii) 由 7.1.1 a) 以及一个\Term{open-subset}{开集}及其补集对应的\Term{cohomology}{上同调}\Term{exact-sequence}{正合列}
\[
 \cdots\longrightarrow H^{i-1}(U,L|U)\xrightarrow{\partial}H_Y^i(X,L)
 \longrightarrow H^i(X,L)\longrightarrow H^i(U,L|U)\longrightarrow\cdots
\]
得到。
\end{proof}

\textbf{注记 7.1.3。} 事实上，由\Term{excision}{切除性}，只要关于 $X$ 的\Term{smoothness}{光滑性}假设以及 $L$ 为\Term{constructible}{可构造}\Term{locally-constant}{局部常值}的假设在 $Y$ 的某个\Term{neighborhood}{邻域}中成立，\Term{canonical}{典范}\Term{isomorphism}{同构} (7.1.2) 仍成立。因此，7.1 (ii) 也可推广到这个范围。

更为一般的表述可见 SGA 2 第 XIV 讲 4.6，不过那里只有借助\Term{resolution-of-singularities}{奇点消解}才能证明，所以目前只适用于\Term{characteristic}{特征} $0$。

\SourcePage{332}
\begin{SGAresult}{推论}{7.2}
设 $X$ 是\Term{dimension}{维数}为 $n$ 的\Term{projective}{射影}、\Term{smooth}{光滑}、\Term{connected}{连通} $k$-\Term{scheme}{概形}，$Y$ 是 $X$ 的一个\Term{subscheme}{子概形}，为 $X$ 的 $d$ 个\Term{hyperplane-section}{超平面截面}之交，$L$ 是一个\Term{constructible}{可构造}\Term{twisted-constant}{扭常值} $\mathbf Z_\ell$-\Term{sheaf}{层}（$\ell\ne\operatorname{car}(k)$）。则：

(i) \Term{inverse-image}{逆像}\Term{homomorphism}{同态}
\[
 \alpha_i:H^i(X,L)\longrightarrow H^i(Y,L|Y)
\]
当 $i\le n-d-1$ 时是\Term{isomorphism}{同构}，当 $i=n-d$ 时是\Term{monomorphism}{单态射}。若再假设 $L$ 是\Term{locally-free}{局部自由}\Term{constructible}{可构造}的，则 $\alpha_{n-d}$ 的\Term{cokernel}{余核}是\Term{torsion-free}{无挠}的。

(ii) 假设 $Y$ 在 $k$ 上\Term{smooth}{光滑}，且\Term{dimension}{维数}为 $n-d$。则 \Term{gysin-homomorphism}{Gysin 同态}
\[
 \beta_i:H^i(Y,(L|Y)(-d))\longrightarrow H^{i+2d}(X,L)
\]
当 $i\ge n-d+1$ 时是\Term{isomorphism}{同构}，当 $i=n-d$ 时是\Term{epimorphism}{满态射}。若再假设 $L$ 是\Term{locally-free}{局部自由}\Term{constructible}{可构造}的，则\Term{epimorphism}{满态射} $\beta_{n-d}$ 是\Term{split-epimorphism}{分裂的}，且其\Term{kernel}{核}是\Term{torsion-free}{无挠}的。
\end{SGAresult}
\begin{proof}
证明 (i)，断言 (ii) 由\Term{poincare-duality}{Poincaré 对偶}得到。在一般情形下，从 7.1 (i) 出发作一次\Term{limit}{取极限}即得结论。当 $L$ 是\Term{locally-free}{局部自由} $\mathbf Z_\ell$-\Term{sheaf}{层}时，对 $D(\mathbf Z_\ell)$ 中的\Term{inverse-image}{逆像}\Term{morphism}{态射}
\[
 \mathbf R\Gamma(X,L)\longrightarrow\mathbf R\Gamma(Y,L|Y)
\]
应用以下\Term{lemma}{引理}，其中假设 (iii) 由 7.1 (i) 得到满足。
\end{proof}

\begin{SGAresult}{\Term{lemma}{引理}}{7.2.1}
设 $V$ 是\Term{discrete-valuation-ring}{离散赋值环}，$k_V$ 是其\Term{residue-field}{剩余域}，$u:K\to L$ 是 $D(V)$ 中的\Term{morphism}{态射}。

假设 $K$ 和 $L$ 的\Term{cohomology}{上同调}均为\Term{finite-type}{有限型}，并以 $M$ 表示 $u$ 的\Term{mapping-cylinder}{映射柱体}（mapping-cylinder）。则对每个 $m\in\mathbf Z$，下列条件\Term{equivalence}{等价}：

(i) 对 $p\le m-1$，有 $H^p(M)=0$，且 $H^m(M)$ 是\Term{torsion-free}{无挠}的。
\SourcePage{333}
(ii) 对每个 $V$-\Term{module}{模} $N$，\Term{morphism}{态射}
\[
 H^p(u\mathbin{\overset{\mathbf L}{\otimes}}\operatorname{id}_N):
 H^p(K\mathbin{\overset{\mathbf L}{\otimes}}N)\longrightarrow
 H^p(L\mathbin{\overset{\mathbf L}{\otimes}}N)
\]
当 $p\le m-1$ 时是\Term{isomorphism}{同构}，当 $p=m$ 时是\Term{monomorphism}{单态射}。

(ii bis) 对每个 $V$-\Term{module}{模} $N$，有
\[
 H^p(M\mathbin{\overset{\mathbf L}{\otimes}}N)=0\qquad(p\le m-1).
\]
(iii) \Term{morphism}{态射}
\[
 H^p(u\mathbin{\overset{\mathbf L}{\otimes}}\operatorname{id}):
 H^p(K\mathbin{\overset{\mathbf L}{\otimes}}k_V)\longrightarrow
 H^p(K\mathbin{\overset{\mathbf L}{\otimes}}k_V)
\]
当 $p\le m-1$ 时是\Term{isomorphism}{同构}，当 $p=m$ 时是\Term{monomorphism}{单态射}。

(iii bis) 对 $p\le m-1$，有
\[
 H^p(M\mathbin{\overset{\mathbf L}{\otimes}}k_V)=0.
\]
此外，这些\Term{equivalence}{等价}条件蕴含：

(i bis) \Term{morphism}{态射}
\[
 H^p(u):H^p(K)\longrightarrow H^p(L)
\]
当 $p\le m-1$ 时是\Term{isomorphism}{同构}，当 $p=m$ 时是\Term{cokernel}{余核}为\Term{free-module}{自由模}的\Term{monomorphism}{单态射}。
\end{SGAresult}
\begin{proof}
下列蕴含图是显然的：
\[
\begin{tikzcd}[column sep=large,row sep=large]
 (\mathrm{ii})\arrow[r,Leftrightarrow]\arrow[d,Rightarrow]&(\mathrm{ii\ bis})
 &(\mathrm{i})\arrow[r,Rightarrow]&(\mathrm{i\ bis})\\
 (\mathrm{iii})\arrow[r,Leftrightarrow]&(\mathrm{iii\ bis})&&
\end{tikzcd}
\]
只需分别证明 $(\mathrm{i})\Rightarrow(\mathrm{ii\ bis})$ 和 $(\mathrm{iii\ bis})\Rightarrow(\mathrm{i})$。由 Künneth \Term{exact-sequence}{正合列}
\[
 0\longrightarrow H^p(M)\otimes_VN\longrightarrow
 H^p(M\mathbin{\overset{\mathbf L}{\otimes}}N)\longrightarrow
 \operatorname{Tor}_1^V(H^{p+1}(M),N)\longrightarrow0
\]
容易得到 $(\mathrm{i})\Rightarrow(\mathrm{ii\ bis})$。

在同一\Term{exact-sequence}{正合列}中取 $N=k_V$，可见若 (iii bis) 成立，则
\[
 H^p(V)\otimes k_V=0\quad(p\le m-1),\qquad
 \operatorname{Tor}_1^V(H^m(M),k_V)=0.
\]
\SourcePage{334}
由 \Term{nakayama-lemma}{Nakayama 引理}及\Term{flatness}{平坦性}的\Term{local-scheme}{局部}判别法（EGA $0_{\mathrm{III}}$ 10.1.3）即得 (i)。
\end{proof}

在上述假设下，知道 $X$ 的 $\mathbf Q_\ell$ 系数\Term{cohomology}{上同调}群，就知道了 $Y$ 的这些群，唯有 $n-1$ 次除外。在这些条件下，后者的计算立即归结为 $Y$ 的\Term{euler-poincare-characteristic}{Euler--Poincaré 示性数}的计算。下面就在稍一般的范围中讨论这一计算。

\begin{SGAresult}{命题}{7.3}
设 $X$ 是\Term{dimension}{维数}为 $n$ 的\Term{proper}{固有}、\Term{smooth}{光滑}、\Term{connected}{连通} $k$-\Term{scheme}{概形}，$E$ 是\Term{rank}{秩}恒为 $r$ 的\Term{locally-free}{局部自由} $\mathscr O_X$-\Term{module}{模}，$Y$ 是 $\mathbf V(E^\vee)$ 的一个与\Term{zero-section}{零截面}\Term{transverse}{横截}的\Term{section}{截面}的零点集。以 $p:X\to k$ 表示\Term{canonical}{典范}\Term{morphism}{态射}，$p_*$ 表示相应的 \Term{gysin-homomorphism}{Gysin 态射}，则在 $H^0(k,\mathbf Q_\ell)=\mathbf Q_\ell$ 中有等式
\[
 \operatorname{EP}(Y)=p_*\bigl(c((\Omega_X^1)^\vee)c(E)^{-1}c_r(E)\bigr).
\]
\end{SGAresult}
\begin{proof}
由 EGA IV 17.13.2，$Y$ 是\Term{smooth}{光滑}的，且纯\Term{dimension}{维数}为 $n-r$。以 $i:Y\to X$ 表示\Term{canonical}{典范}\Term{immersion}{浸入}，则\Term{conormal-sheaf}{余法层} $\mathscr N_i$ 是 $\mathbf V(E^\vee)$ 沿\Term{zero-section}{零截面}的\Term{conormal-sheaf}{余法层} $E^\vee$ 经 $i$ 的\Term{inverse-image}{逆像}。因此有\Term{exact-sequence}{正合列}
\[
 0\longrightarrow i^*(E^\vee)\longrightarrow i^*(\Omega_X^1)
 \longrightarrow\Omega_Y^1\longrightarrow0,\tag{7.3.1}
\]
由此立即得到等式
\[
 c((\Omega_Y^1)^\vee)=i^*\bigl(c((\Omega_X^1)^\vee)c(E)^{-1}\bigr).\tag{7.3.2}
\]
再由 (3.10) 得到 $\operatorname{EP}(Y)$ 的如下表达式：
\[
 \operatorname{EP}(Y)=p_*i_*i^*\bigl(c((\Omega_X^1)^\vee)c(E)^{-1}\bigr).
\]
利用\Term{morphism}{态射} $i$ 的\Term{projection-formula}{投影公式}，可写为
\[
 \operatorname{EP}(Y)=p_*\bigl(c((\Omega_X^1)^\vee)c(E)^{-1}i_*(1_Y)\bigr).
\]
由 (4.7) 即得命题中的公式。
\end{proof}

\SourcePage{335}
\begin{SGAresult}{推论}{7.4}
设 $X$ 是\Term{dimension}{维数}为 $n$ 的\Term{proper}{固有}、\Term{smooth}{光滑}、\Term{connected}{连通} $k$-\Term{scheme}{概形}，$Y$ 是 $X$ 的一个\Term{dimension}{维数}为 $n-1$ 的\Term{smooth}{光滑}\Term{closed-subscheme}{闭子概形}，因而由某个\Term{invertible-module}{可逆 $\mathscr O_X$-模} $L$ 的一个\Term{section}{截面}定义。以 $p$ 表示 $X$ 到 $k$ 的\Term{canonical}{典范}\Term{morphism}{态射}，$p_*$ 表示相应的 \Term{gysin-homomorphism}{Gysin 态射}，则在 $H^0(k,\mathbf Q_\ell)=\mathbf Q_\ell$ 中有等式
\[
 \operatorname{EP}(X)-\operatorname{EP}(Y)
 =p_*\bigl(c((\Omega_X^1)^\vee)c(L)^{-1}\bigr).\tag{7.4.1}
\]
此外，
\[
 b^{n-1}(Y)=2b^{n-1}(X)-b^n(X)
 +(-1)^np_*\bigl(c((\Omega_X^1)^\vee)c(L)^{-1}\bigr).\tag{7.4.2}
\]
\end{SGAresult}
\begin{proof}
考虑显然的等式
\[
 1-c(L)^{-1}c_1(L)=c(L)^{-1},
\]
由 7.3 和 4.2 即得等式 (7.4.1)。由 (7.2) 可见
\[
 \begin{split}
 b^{n-1}(Y)={}&b^{n-1}(X)-b^n(X)+b^{n+1}(X)\\
 &+(-1)^n\bigl(\operatorname{EP}(X)-\operatorname{EP}(Y)\bigr).
 \end{split}
\]
但由\Term{poincare-duality}{Poincaré 对偶}，$b^{n-1}(X)=b^{n+1}(X)$，故有 (7.4.2)。
\end{proof}

\begin{SGAresult}{推论}{7.5}
设 $X$ 是 $\mathbf P_k^{n+r}$ 的一个\Term{smooth}{光滑}、\Term{connected}{连通}、\Term{dimension}{维数}为 $n$ 的\Term{subscheme}{子概形}，是 $\mathbf P_k^{n+r}$ 中 $r$ 个\Term{divisor}{除子} $X_1,X_2,\ldots,X_r$ 的\Term{global-complete-intersection}{整体完全交}。对每个整数 $i\in(1,r)$，以 $d_i$ 表示 $X_i$ 的\Term{degree}{次数}。则：

(i) \Term{cohomology}{上同调}群 $H^p(X,\mathbf Z_\ell)$ 是\Term{finite-type}{有限型}\Term{free-module}{自由 $\mathbf Z_\ell$-模}。

(ii) 若整数 $p\le2n$ 且 $p\ne n$，则
\[
 H^p(X,\mathbf Z_\ell)=
 \begin{cases}
 0,&p\text{ 为奇数},\\
 \mathbf Z_\ell,&p\text{ 为偶数}.
 \end{cases}
\]
(iii) $X$ 的\Term{euler-poincare-characteristic}{Euler--Poincaré 示性数}由下式给出：
\[
 \operatorname{EP}(X)=d_1\cdots d_r
 \sum_{\substack{0\le i\le n\\j_1+\cdots+j_r=n-i}}
 (-1)^{n-i}\binom{n+r+1}{i}d_1^{j_1}\cdots d_r^{j_r}.
\]
\end{SGAresult}
\begin{proof}
断言 (i) 立即由 (7.2) 得到。考虑\Term{projective-space}{射影空间}的\Term{cohomology}{上同调}（第 2 节），(ii) 也同样得到。证明 (iii)。以 $P$ 表示周围的\Term{projective-space}{射影空间}，$p:P\to k$ 表示\Term{canonical-projection}{典范投影}，$p_*$ 表示相应的 \Term{gysin-homomorphism}{Gysin 态射}，
\SourcePage{336}
$\xi$ 表示 $\mathscr O_P(1)$ 的类。对每个 $i\in(1,r)$，\Term{scheme}{概形} $X_i$ 由某个\Term{invertible-module}{可逆 $\mathscr O_P$-模} $L_i$ 的一个\Term{section}{截面}定义，因此 $X$ 是 $\mathbf V(L_1^\vee\oplus\cdots\oplus L_r^\vee)$ 的一个与\Term{zero-section}{零截面}\Term{transverse}{横截}的\Term{section}{截面}的零点集。另一方面，已知有\Term{exact-sequence}{正合列}
\[
 0\longrightarrow\Omega_P^1\otimes\mathscr O_P(1)
 \longrightarrow\mathscr O_P^{n+r+1}\longrightarrow\mathscr O_P(1)\longrightarrow0,
\]
故
\[
 c((\Omega_P^1)^\vee)=(1+\xi)^{n+r+1}.
\]
由 (7.3) 得到 $\operatorname{EP}(X)$ 的如下表达式：
\[
 \begin{split}
 \operatorname{EP}(X)
 &=p_*\bigl(c((\Omega_P^1)^\vee)c(L_1^\vee\oplus\cdots\oplus L_r^\vee)^{-1}
 c_1(L_1^\vee)\cdots c_1(L_r^\vee)\bigr)\\
 &=p_*\bigl(d_1\cdots d_r\xi^r(1+\xi)^{n+r+1}
 (1+d_1\xi)^{-1}\cdots(1+d_r\xi)^{-1}\bigr).
 \end{split}
\]
利用 $p_*(\xi^{n+r})=1$（2.2.2 c)），容易得出命题中的公式。
\end{proof}

\textbf{注记 7.6。} 也可利用 (4.11)，并转写 Hirzebruch 借助 \Term{riemann-roch-theorem}{Riemann--Roch 定理}所作的计算（[3]，第 160--161 页），得到\Term{euler-poincare-characteristic}{Euler--Poincaré 示性数}的表达式。关于更完整的计算以及\Term{complete-intersection}{完全交}的 Hodge \Term{cohomology}{上同调}的确定，可参阅 SGA 7 第 XI 讲。

\textbf{7.7。} 在超越情形中，对\Term{constant-coefficients}{常系数}还有另一个众所周知的 Lefschetz \Term{theorem}{定理}（[4]）。与此类比，便自然提出如下猜想。\footnote{正如 Deligne 所指出的，若不对 $L$ 加上补充假设，(7.7) 中的断言是错误的。不过，当 $L$ 还满足“\Term{pure}{纯}”的条件时，Deligne 已经证明了这些断言（《La conjecture de Weil II》）。}

若 $X$ 是\Term{dimension}{维数}为 $n$ 的\Term{projective}{射影}、\Term{smooth}{光滑}、\Term{connected}{连通} $k$-\Term{scheme}{概形}，$\xi$ 是某个\Term{very-ample}{很丰沛}的\Term{invertible-module}{可逆 $\mathscr O_X$-模}在 $H^2(X,\mathbf Z_\ell(1))$ 中的类，则对每个\Term{twisted-constant}{扭常值} $\mathbf Q_\ell$-\Term{sheaf}{层} $L$ 和每个整数 $p\le n$，乘以 $\xi^{n-p}$ 的映射
\[
 H^p(X,L)\longrightarrow H^{2n-p}(X,L(n-p))
\]
是\Term{isomorphism}{同构}。

\SourcePage{337}
这一猜想断言有如下推论，它们在\Term{characteristic}{特征}非零的情形下也尚未得到证明。

设 $X$ 是\Term{dimension}{维数}为 $n$ 的\Term{projective}{射影}、\Term{smooth}{光滑}、\Term{connected}{连通} $k$-\Term{scheme}{概形}，$Y$ 是 $X$ 的一个\Term{smooth}{光滑}\Term{hyperplane-section}{超平面截面}，$L$ 是 $X$ 上的一个\Term{twisted-constant}{扭常值} $\mathbf Q_\ell$-\Term{sheaf}{层}。

(i) \Term{inverse-image}{逆像}\Term{morphism}{态射}
\[
 H^p(X,L)\longrightarrow H^p(Y,L|Y)
\]
当 $p\ge n$ 时是\Term{epimorphism}{满态射}。

(ii) \Term{gysin-homomorphism}{Gysin 态射}
\[
 H^p(Y,(L|Y)(-1))\longrightarrow H^{p+2}(X,L)
\]
当 $p\le n-2$ 时是\Term{monomorphism}{单态射}。

\begin{proof}
设 $j:X\hookrightarrow\mathbf P^r$ 是\Term{closed-immersion}{闭浸入}，$\xi_X$（相应地，$\xi_Y$）是 $\mathscr O_X(1)$（相应地，$\mathscr O_Y(1)$）在 $H^2(X,\mathbf Q_\ell(1))$（相应地，$H^2(Y,\mathbf Q_\ell(1))$）中的类。以 $\alpha:Y\hookrightarrow X$ 表示\Term{canonical}{典范}\Term{injection}{单射}。对 $q\le n-1$，下图\Term{commutative-diagram}{交换}：
\[
\begin{tikzcd}[column sep=large,row sep=large]
 H^q(X,L)\arrow[r,"\alpha^*"]\arrow[d,"{\cup\,\xi_X^{n-1-q}}"']
 &H^q(Y,L|Y)\arrow[d,"{\cup\,\xi_Y^{n-1-q}}"]\\
 H^{2n-2-q}(X,L(n-1-q))\arrow[r,"\alpha^*"']
 &H^{2n-2-q}(Y,(L|Y)(n-1-q)).
\end{tikzcd}
\]
右列（按猜想）是\Term{isomorphism}{同构}。置 $p=2n-2-q$；当 $q\le n-2$，即 $p\ge n$ 时，上行是\Term{isomorphism}{同构}（7.2），从而得到断言 (i)。断言 (ii) 的证明类似，这次对 $p\le n-1$ 考虑下图：
\SourcePage{338}
\[
\begin{tikzcd}[column sep=large,row sep=large]
 H^p(Y,(L|Y)(-1))\arrow[r,"\mathrm{Gysin}"]\arrow[d,"{\cup\,\xi_Y^{n-1-p}}"']
 &H^{p+2}(X,L)\arrow[d,"{\cup\,\xi_X^{n-1-p}}"]\\
 H^{2n-2-p}(Y,(L|Y)(n-2-p))\arrow[r,"\mathrm{Gysin}"']
 &H^{2n-p}(X,L(n-1-p)).
\end{tikzcd}
\]
其中当 $2n-2-p\ge n$，即 $p\le n-2$ 时，下行是\Term{isomorphism}{同构}（7.2 (ii)）。
\end{proof}

关于这些猜想及其推论的讨论，可参阅 Kleiman 的报告《Algebraic cycles and the Weil conjectures》，载于《Dix exposés sur la cohomologie des schémas》。
